\documentclass[11pt,a4paper]{article}

\usepackage[margin=2.5cm]{geometry}
\usepackage{graphicx}
\usepackage{booktabs}
\usepackage{array}
\usepackage{tabularx}
\usepackage{amsmath}
\usepackage{amssymb}
\usepackage{hyperref}
\usepackage{enumitem}
\usepackage{xcolor}

\hypersetup{
    colorlinks=true,
    linkcolor=blue,
    urlcolor=blue,
    citecolor=blue
}

\title{AMD RAG Vector Search Gap Demo\\[0.4em]
\large Project Plan (3--4 Engineers, 2--3 Weeks)}
\author{Konstantin Kolchin}
\date{\today}

\begin{document}
\maketitle

\begin{abstract}
This document proposes a focused engineering sprint to demonstrate the
\emph{integration gap} between NVIDIA-centric GPU vector search in production
RAG stacks and the AMD ROCm ecosystem.
Milvus~\cite{milvus-home,milvus-osdi21} exposes GPU-accelerated indexing and
query through RAPIDS cuVS~\cite{cuvs-docs,rapids-home}, which requires CUDA;
AMD offers hipVS~\cite{hipvs-docs,hipvs-github}---a ROCm-DS port of cuVS---but
no first-class Milvus integration exists (GitHub issue~\#5654 remains
unresolved~\cite{milvus-rocm-issue}).
Over two to three weeks, a team of three to four engineers will produce a gap
analysis, a working hipVS retrieval demo embedded in a minimal RAG pipeline,
recall/latency benchmarks against CPU baselines, and a Milvus integration spike
with effort estimates.
The sprint deliberately excludes full-stack ports (Milvus, RAPIDS, Audacity,
OpenVINO) and instead delivers evidence and a reproducible demo suitable for
stakeholder review.
\end{abstract}

\section{Background and motivation}
\label{sec:background}

\subsection{Retrieval-augmented generation (RAG)}

RAG augments large language models with external knowledge by retrieving
relevant documents at query time~\cite{rag-lewis2020,langchain-rag,llamaindex-rag}.
A typical pipeline embeds corpus chunks into dense vectors
(e.g.\ via Sentence-Transformers~\cite{sentence-transformers} or commercial
APIs~\cite{openai-embeddings}), stores them in a vector database, and at query
time retrieves the top-$k$ nearest neighbors before passing context to the LLM.
Embedding quality is commonly evaluated on benchmarks such as
MTEB~\cite{mteb}.

Vector search---finding approximate nearest neighbors (ANN) in high-dimensional
space at scale---is the performance-critical step once corpora exceed in-memory
CPU capacity.
Libraries and databases in this space include FAISS~\cite{faiss-github,faiss-arxiv},
Milvus~\cite{milvus-home,milvus-github}, Chroma~\cite{chromadb}, Qdrant~\cite{qdrant},
Weaviate~\cite{weaviate}, and managed offerings such as Pinecone~\cite{pinecone-arch}.

\subsection{GPU-accelerated vector search on NVIDIA}

NVIDIA RAPIDS cuVS~\cite{cuvs-docs,cuvs-github,rapids-home,rapids-github}
provides GPU implementations of ANN algorithms including IVF-PQ~\cite{ivfpq-jegou},
HNSW-related methods~\cite{hnsw-paper}, and CAGRA~\cite{nvidia-cagra}.
Milvus integrates cuVS for GPU index build and search~\cite{milvus-gpu-index,milvus-gpu-cagra},
making CUDA GPUs the de facto accelerator path for production Milvus deployments
that require sub-millisecond query latency on large collections.

\subsection{AMD ROCm and hipVS}

AMD's HIP~\cite{hip-docs,hip-porting,hipify-device-api} provides a CUDA-like
programming model on ROCm~\cite{rocm-prog-guide}.
The ROCm Data Science (ROCm-DS) initiative ports key RAPIDS components to HIP;
hipVS~\cite{hipvs-docs,hipvs-whatis,hipvs-github,rocm-ds-overview} is the
vector-search counterpart to cuVS, with hipRAFT~\cite{hipraft-github} backing
low-level primitives.
AMD documents build instructions~\cite{hipvs-build} and reports performance
characteristics in ROCm blog posts~\cite{hipvs-blog}.

Prior work in this repository~\cite{cuda-to-hip-repo} demonstrated that
CUDA-to-HIP kernel ports can yield significant wins on AMD hardware when
architecture-specific tuning (LDS layout, divergence reduction) is applied;
the same engineering mindset applies to integrating hipVS into application stacks
rather than assuming API parity alone guarantees performance.

\subsection{The gap}

Milvus GPU support is CUDA-only~\cite{milvus-gpu-index,milvus-prerequisite-gpu}.
Official prerequisites list NVIDIA compute capability, drivers, and the
\texttt{nvidia-device-plugin} on Kubernetes~\cite{milvus-prerequisite-gpu};
GPU index types such as \texttt{GPU\_CAGRA} and \texttt{GPU\_IVF\_PQ} are
contributed through NVIDIA RAPIDS/cuVS~\cite{milvus-gpu-cagra,milvus-gpu-index}.
Community requests for ROCm/AMD support~\cite{milvus-rocm-issue} have not
produced a merged integration path.
Meanwhile hipVS exists and is actively maintained~\cite{hipvs-github}, but
end users cannot select ``Milvus + AMD GPU'' as they can ``Milvus + NVIDIA GPU.''
This sprint makes that gap \emph{visible and measurable}: not merely a feature
request, but a documented user journey, benchmark data, and integration options.

\paragraph{Illustrative journeys from community reports.}
The following vignettes are drawn from public Milvus GitHub issues and
discussions; they are representative of recurring reports, not isolated bugs.

\medskip
\noindent\textbf{On NVIDIA (straightforward path).}
A team building a RAG service on Linux with an NVIDIA A100 or T4 follows
documented steps: install the GPU-enabled Milvus standalone or cluster compose
file~\cite{milvus-gpu-index,milvus-230-beta}, create a collection, and build a
\texttt{GPU\_CAGRA} or \texttt{GPU\_IVF\_PQ} index.
Milvus documentation states that GPU indexing can improve throughput by up to
$100\times$ on large batches relative to CPU~\cite{milvus-gpu-index}.
Embedding generation (Sentence-Transformers or an API) plus Milvus GPU retrieval
is a well-trodden path in tutorials and production write-ups~\cite{langchain-rag,milvus-home}.

\medskip
\noindent\textbf{On AMD or ROCm-class hardware (blocked at the database layer).}
In June~2021, GitHub user \texttt{boniek83} opened a feature request titled
``ROCm support,'' stating plainly that ``Milvus only works with NVidia
gpus'' and asking for a ROCm build~\cite{milvus-rocm-issue}.
The issue was moved to the 2.x backlog and closed as stale three months later
without implementation---a pattern repeated across subsequent years.

In March~2023, newcomer \texttt{victorsoda} asked on Milvus Discussions how to
run standalone Milvus \emph{on their GPUs} for indexing and similarity
search, noting that their chip is ``similar to GPU but different'' and asking
whether they could contribute support~\cite{milvus-discussion-23145}.
Milvus maintainer \texttt{yhmo} replied that GPU support exists from v2.3.0-beta
onward but ``only supports Nvidia cards and requires NVIDIA Driver
installed''~\cite{milvus-discussion-23145,milvus-230-beta}.
There is no documented path for ROCm or AMD Radeon in that thread---only the
NVIDIA prerequisite chain.

A structurally similar report appears for Hygon DCU accelerators (ROCm/HIP
ecosystem, AMD-derived architecture).
In October~2025, \texttt{deqinxirang} asked whether Milvus supports the Hygon
K100\_AI DCU and \emph{what steps are required to add support}~\cite{milvus-issue-dcu}.
Core maintainer \texttt{xiaofan-luan} answered: ``don't have plan right now for
any other device rather than GPU'' (meaning NVIDIA GPU in practice), ``But any
contributions are welcomed''~\cite{milvus-issue-dcu}.
The author followed up weeks later---``Does the latest version support Hygon DCU
now? How should one adapt?''---and received no integration guide before the
issue was closed stale~\cite{milvus-issue-dcu,milvus-accelerator-issue}.

Broader accelerator requests tell the same story at the architectural level.
Issue~\#21550 asks about Kunlun, Hygon DCU, and other AI cards~\cite{milvus-accelerator-issue};
issue~\#33885 notes that Knowhere depends on NVIDIA RAFT/cuVS, which ``only supports
CUDA,'' so non-CUDA GPUs require rewriting ANN kernels for another stack~\cite{milvus-huawei-issue}.
Milvus engineers describe a roadmap for flexible third-party devices but defer
implementation to community contributors~\cite{milvus-accelerator-issue,milvus-huawei-issue}.

\medskip
\noindent\textbf{Composite RAG scenario (synthesized from the above).}
Consider a developer who already runs embeddings on an AMD workstation
(Radeon RX~6800 XT or an institutional Hygon DCU node under ROCm).
They read Milvus RAG tutorials~\cite{langchain-rag,rag-lewis2020}, install Milvus,
and expect to enable \texttt{GPU\_CAGRA} as in the NVIDIA examples~\cite{milvus-gpu-index}.
Instead they discover that GPU indexes never register: Milvus looks for CUDA and
NVIDIA drivers~\cite{milvus-prerequisite-gpu,milvus-discussion-23145}.
CPU-only Milvus remains usable but falls short when documentation promises
order-of-magnitude GPU gains~\cite{milvus-gpu-index}.
AMD's hipVS~\cite{hipvs-docs,hipvs-github} offers the missing ANN layer on ROCm,
yet there is no \texttt{pymilvus} switch to select it---the user must leave the
Milvus GPU path entirely and wire hipVS (or FAISS on CPU) into a custom RAG
script~\cite{faiss-github,hipvs-blog}.
That integration friction---not the absence of AMD vector-search code---is the
gap this sprint documents and demos.

\section{Project goals}
\label{sec:goals}

\begin{enumerate}[label=\textbf{G\arabic*.}, leftmargin=*]
  \item \textbf{Gap analysis document} mapping CUDA/NVIDIA vs.\ ROCm/AMD capabilities
        across vector search, Milvus, and RAG tooling, with a prioritized integration
        matrix.
  \item \textbf{hipVS demo} building IVF and/or CAGRA indexes on AMD hardware,
        reporting recall@$k$ and query latency vs.\ CPU/FAISS baselines.
  \item \textbf{Minimal RAG pipeline} wiring embeddings $\rightarrow$ hipVS retrieval
        $\rightarrow$ sample queries (CLI or notebook), proving end-to-end feasibility
        outside Milvus.
  \item \textbf{Milvus integration spike} identifying where GPU code paths branch on
        CUDA, sketching hipVS adapter options, and estimating effort for a production-grade
        contribution.
\end{enumerate}

\section{Scope}
\label{sec:scope}

\subsection{In scope}

\begin{itemize}[leftmargin=*]
  \item Literature and product survey (Milvus, cuVS, hipVS, FAISS, RAG frameworks).
  \item hipVS build and validation on a ROCm machine (target: \texttt{gfx1030} or
        newer RDNA/CDNA as available).
  \item Benchmark suite: index build time, queries/sec, recall@$k$ on a public or
        synthetic embedding dataset.
  \item Demo RAG script (Python): chunk corpus, embed, index with hipVS, retrieve,
        optional LLM call stub.
  \item Milvus source review: GPU index factory, cuVS bindings, configuration surface.
  \item Written deliverables: gap report, demo README, integration proposal.
\end{itemize}

\subsection{Out of scope (non-goals)}

\begin{itemize}[leftmargin=*]
  \item Full Milvus fork with merged ROCm support (spike only).
  \item Port of entire RAPIDS stack to ROCm beyond what hipVS/hipRAFT already provide.
  \item Audacity~\cite{audacity-home} or OpenVINO Audacity
        plugins~\cite{openvino-audacity} (orthogonal ONNX/ROCm paths exist
        separately~\cite{onnxruntime-rocm}).
  \item Production hardening: HA Milvus cluster, auth, multi-tenant ops.
  \item Training or fine-tuning embedding models.
\end{itemize}

\section{Workstreams and staffing}
\label{sec:workstreams}

Recommended team: \textbf{3--4 engineers} over \textbf{2--3 weeks}.
Table~\ref{tab:workstreams} assigns primary ownership; all streams share weekly
sync and a common benchmark harness.

\begin{table}[htbp]
  \centering
  \caption{Workstreams, owners, and primary deliverables.}
  \label{tab:workstreams}
  \begin{tabularx}{\textwidth}{@{}l X X@{}}
    \toprule
    \textbf{Stream} & \textbf{Focus} & \textbf{Deliverable} \\
    \midrule
    W1 -- Gap analysis &
    Survey CUDA vs.\ AMD matrix; user journeys; risk register &
    Gap analysis PDF (15--25 pages) \\
    W2 -- hipVS engineer &
    Build hipVS~\cite{hipvs-build}; IVF/CAGRA; recall/latency &
    Benchmark tables + reproducible scripts \\
    W3 -- RAG pipeline &
    Embeddings~\cite{sentence-transformers}; retrieval demo; docs &
    \texttt{demo/} repo + runbook \\
    W4 -- Milvus spike &
    Trace GPU paths~\cite{milvus-gpu-index}; adapter design; estimate &
    Integration memo + optional PoC patch \\
    \bottomrule
  \end{tabularx}
\end{table}

\subsection{W1: Gap analysis (lead: systems architect)}

\textbf{Week 1:}
\begin{itemize}[leftmargin=*]
  \item Inventory GPU vector search APIs: cuVS~\cite{cuvs-docs} vs.\
        hipVS~\cite{hipvs-docs}.
  \item Document Milvus GPU architecture~\cite{milvus-gpu-cagra,milvus-osdi21} and
        ROCm issue history~\cite{milvus-rocm-issue}.
  \item Compare alternative stacks (FAISS~\cite{faiss-arxiv}, Milvus Lite~\cite{milvus-lite})
        for AMD-only deployments.
\end{itemize}

\textbf{Week 2--3:}
\begin{itemize}[leftmargin=*]
  \item Produce integration matrix (component $\times$ CUDA $\times$ ROCm $\times$ effort).
  \item Write executive summary for Alexander/stakeholders.
  \item Cross-link W2/W3 benchmark numbers into the report.
\end{itemize}

\subsection{W2: hipVS engineer (lead: GPU/ROCm specialist)}

\textbf{Week 1:}
\begin{itemize}[leftmargin=*]
  \item Clone and build ROCm-DS/hipVS~\cite{hipvs-github} on target AMD machine.
  \item Run upstream examples; verify HIP/ROCm version compatibility~\cite{hip-faq}.
  \item Select datasets (e.g.\ 768-dim sentence embeddings, $10^5$--$10^6$ vectors).
\end{itemize}

\textbf{Week 2:}
\begin{itemize}[leftmargin=*]
  \item Benchmark IVF-PQ and CAGRA (if available) vs.\ CPU FAISS~\cite{faiss-github}.
  \item Measure build time, batch query latency, recall@1/10/100.
  \item Profile bottlenecks; note hipRAFT dependency behavior~\cite{hipraft-github}.
\end{itemize}

\textbf{Week 3:}
\begin{itemize}[leftmargin=*]
  \item Harden scripts for reproducibility; document hardware (\texttt{gfx*} arch).
  \item Support W3 integration; optional comparison with cuVS docs baseline~\cite{nvidia-cagra}.
\end{itemize}

\subsection{W3: RAG pipeline (lead: ML/application engineer)}

\textbf{Week 1:}
\begin{itemize}[leftmargin=*]
  \item Define demo corpus (technical docs or Wikipedia subset).
  \item Pipeline: chunk $\rightarrow$ embed (Sentence-Transformers~\cite{sentence-transformers})
        $\rightarrow$ numpy/torch tensors.
\end{itemize}

\textbf{Week 2:}
\begin{itemize}[leftmargin=*]
  \item Wire hipVS index from W2; expose \texttt{search(query, k)}.
  \item Implement query loop following RAG patterns~\cite{langchain-rag,rag-lewis2020}.
  \item CLI or Jupyter demo with 5--10 example questions.
\end{itemize}

\textbf{Week 3:}
\begin{itemize}[leftmargin=*]
  \item README with prerequisites (ROCm, hipVS, Python deps).
  \item Optional: stub LLM response (no API key required for core demo).
\end{itemize}

\subsection{W4: Milvus integration spike (lead: backend engineer)}

\textbf{Week 1:}
\begin{itemize}[leftmargin=*]
  \item Clone milvus-io/milvus~\cite{milvus-github}; map GPU index code paths.
  \item Identify cuVS call sites and CUDA assumptions.
\end{itemize}

\textbf{Week 2:}
\begin{itemize}[leftmargin=*]
  \item Design options: (A)~hipVS backend plugin, (B)~sidecar index service,
        (C)~Milvus CPU + external hipVS retrieval.
  \item Estimate LOC, test burden, upstream contribution path.
\end{itemize}

\textbf{Week 3:}
\begin{itemize}[leftmargin=*]
  \item Write integration memo (recommended path + risks).
  \item Optional: minimal patch or fork branch demonstrating index factory hook.
\end{itemize}

\section{Timeline}
\label{sec:timeline}

\begin{table}[htbp]
  \centering
  \caption{Three-week sprint calendar (adjustable to two weeks by deferring W4 PoC).}
  \label{tab:timeline}
  \begin{tabular}{@{}llll@{}}
    \toprule
    \textbf{Week} & \textbf{W1 Gap} & \textbf{W2 hipVS} & \textbf{W3 RAG / W4 Milvus} \\
    \midrule
    1 & Survey + outline & Build + smoke tests & Corpus + Milvus trace \\
    2 & Draft report & Benchmarks & hipVS wire + adapter design \\
    3 & Final report & Repro scripts & Demo polish + integration memo \\
    \bottomrule
  \end{tabular}
\end{table}

\textbf{Milestone M1 (end of week 1):} hipVS builds; Milvus GPU map drafted.
\textbf{M2 (week 2):} First benchmark numbers; RAG retrieves correct chunks on sample queries.
\textbf{M3 (week 3):} All deliverables tagged; stakeholder demo ready.

\section{Success criteria}
\label{sec:success}

\begin{enumerate}[label=\textbf{S\arabic*.}, leftmargin=*]
  \item Gap report cites primary sources (Milvus docs, hipVS docs, issue~\#5654) and
        includes a clear ``what works today on AMD'' vs.\ ``what Milvus users expect'' table.
  \item hipVS demo achieves recall within 5\% of CPU FAISS baseline at comparable $k$ on
        at least one dataset.
  \item RAG demo returns relevant chunks for $\geq$80\% of hand-crafted test queries
        (qualitative review).
  \item Milvus spike delivers a ranked list of integration options with person-week estimates.
  \item All artifacts in a single repository with CI-friendly build/run instructions.
\end{enumerate}

\section{Risks and mitigations}
\label{sec:risks}

\begin{table}[htbp]
  \centering
  \caption{Risk register.}
  \label{tab:risks}
  \begin{tabularx}{\textwidth}{@{}%
    >{\raggedright\arraybackslash}p{0.32\textwidth}%
    >{\raggedright\arraybackslash}p{0.14\textwidth}%
    >{\raggedright\arraybackslash}X@{}}
    \toprule
    \textbf{Risk} & \textbf{Impact} & \textbf{Mitigation} \\
    \midrule
    hipVS build fails on target GPU &
    Blocks W2/W3 &
    Early week-1 build; fallback to hipVS Docker/CI runner \\
    Algorithm parity gaps (CAGRA vs IVF) &
    Weak demo &
    Document supported index types~\cite{hipvs-whatis}; use best available \\
    Milvus codebase complexity &
    W4 overrun &
    Time-box spike; memo over merge-ready PR \\
    Embedding model size on CPU &
    Slow demo prep &
    Cache embeddings; use small model \\
    Stakeholder expects full Milvus port &
    Scope creep &
    Refer to non-goals (Section~\ref{sec:scope}) \\
    \bottomrule
  \end{tabularx}
\end{table}

\section{Relationship to prior HIP kernel work}
\label{sec:prior}

The isolated-ones HIP benchmark in \texttt{cuda-to-hip}~\cite{cuda-to-hip-repo}
showed that portable CUDA subsets (\texttt{\_\_shared\_\_}, \texttt{\_\_shfl},
\texttt{hipEvent}) suffice for custom kernels, but micro-architecture tuning
(LDS bank padding, divergence) determines winners on \texttt{gfx1030}.
hipVS abstracts ANN algorithms behind a library API~\cite{hipvs-docs}; this sprint
shifts focus from hand-written kernels to \emph{stack integration}---yet the
same validation discipline applies: measure on AMD hardware, do not assume
CUDA-era defaults transfer unchanged~\cite{hip-porting,hipify-device-api}.

\section{Deliverables checklist}
\label{sec:deliverables}

\begin{itemize}[leftmargin=*]
  \item[\checkmark] \textbf{D1:} Gap analysis report (PDF).
  \item[\checkmark] \textbf{D2:} \texttt{benchmarks/} --- hipVS vs.\ FAISS scripts and results CSV.
  \item[\checkmark] \textbf{D3:} \texttt{demo/rag\_hipvs/} --- minimal RAG pipeline.
  \item[\checkmark] \textbf{D4:} \textbf{Milvus integration memo} (PDF or markdown).
  \item[\checkmark] \textbf{D5:} 30-minute stakeholder walkthrough (slides optional).
\end{itemize}

\section{References}

\bibliographystyle{plainurl}
\bibliography{amd_rag_gap_plan}

\end{document}
